\section{Introduction}
\label{sec:introduction}

Exact feasibility questions for neural networks sit at the intersection of
real algebraic geometry, computational complexity, and formal verification.
Once approximation and numerical optimization are removed, a finite
feed-forward ReLU network with unknown real weights and rational data defines a
finite system of polynomial equalities and inequalities.  The natural decision
question is whether that system has a real solution.  Closely related neural
network training problems are known to be complete for the complexity class
\(\ExistsR\)~\cite{AbrahamsenKleistMiltzow2021}; this class contains NP and is
contained in PSPACE~\cite{Canny1988,SchaeferCardinalMiltzow2024}.

This paper investigates whether the repeated-sample structure of exact neural
training admits an exact separator algorithm.  A single parameter vector
\(\theta\) is shared across the dataset, while each sample has its own
preactivations, activations, and auxiliary variables.  At the level of
polynomial equations this yields a block-angular Jacobian.  At the level of
commutative algebra it yields an iterated tensor product over the shared
parameter algebra.  These facts motivate \emph{Cotangent Message Passing}
(CMP): eliminate a local block from the row space of a linearized constraint
system and retain only the induced subspace on a separator.

The distinction between a linearized relation and an exact existential
projection is decisive.  The cotangent module records first-order behavior.
The decision problem asks whether the semialgebraic set is empty.  Neither the
universal property of K\"ahler differentials nor a pointwise Jacobian identity
identifies these two objects.  A proof of an exact CMP decider therefore needs
an additional semantic invariant:
\begin{equation}
  s\in \mathsf{Msg}_v
  \quad\Longleftrightarrow\quad
  \exists x_{\operatorname{int}(T_v)}\;
  F_{T_v}(x,s).
  \label{eq:semantic-message-intro}
\end{equation}
It must also prove that the representation of \(\mathsf{Msg}_v\), its merge
operation, and its projection operation remain polynomially bounded.

\subsection{Contributions}

The present draft makes the following claims, each with an explicit proof
status.

\begin{enumerate}
  \item We define \(\ExactNNTP\) as a Boolean feasibility language with
  rational input and unrestricted real witnesses.  ``Training'' describes the
  origin of the constraints; the output is a Boolean, not a weight vector.

  \item We prove an exact complementarity description of a ReLU gate and its
  existential square-slack polynomialization.  The description includes the
  switching point \(z=a=0\).

  \item We derive the shared-base tensor-product description of the equality
  coordinate ring and the direct-sum decomposition of relative K\"ahler
  differentials after base change.  We then display the induced block-angular
  Jacobian.

  \item We define the intrinsic first-order CMP message
  \(
    \cE(A\mathbin{\to}S)=J_{AS}(\kernel J_A)
  \)
  and prove exact separator factorization for the associated linear row-space
  problem.  This theorem is independent of a null-space basis.

  \item We prove that the first-order message is not a complete invariant of
  exact ReLU feasibility.  A one-ReLU family with positive target \(y\) has
  exact parameter projection \(\{y\}\), but its evaluated scalar CMP message
  equals \(\R\) for every \(y>0\).  Thus targets one and two yield identical
  messages and different exact feasibility at \(\theta=1\).

  \item We isolate the remaining obligations for the proposed complexity
  collapse: a polynomially represented zeroth-order message, a correct global
  decider, polynomial Turing-model bit complexity, and a reduction landing in
  the exact language defined here.  We formalize the final class-theoretic
  implication conditionally.
\end{enumerate}

\subsection{Status of the proposed collapse}

The intended implication has the form
\begin{equation}
  \ExistsR \leq_p \ExactNNTP \in P
  \quad\Longrightarrow\quad
  P=NP=\ExistsR.
  \label{eq:collapse-chain}
\end{equation}
The implication in \cref{eq:collapse-chain} is standard.  The right-hand
premise is not supplied by the first-order CMP theorem.  Indeed,
\cref{sec:global-exactness} disproves the required completeness statement for
a decoder that receives only the existing cotangent message subspace.  The
collapse is therefore recorded as a target conditional theorem, not as an
unconditional result.

This separation is more than a matter of exposition.  The counterexample is
inside the stated problem class, uses a single ReLU, and occurs at a positive
satisfying point.  It cannot be repaired by excluding ReLU switching
singularities.  Any successful exact message must preserve zeroth-order target
constants and global component information that differentiation can erase.

\subsection{Organization}

\Cref{sec:preliminaries,sec:problem,sec:relu} define the computational and
algebraic setting.  \Cref{sec:coordinate-ring,sec:kahler,sec:block-jacobian}
establish the shared-base structure.  \Cref{sec:topology,sec:cmp} develop the
separator decomposition and the proved linear CMP theorem.
\Cref{sec:global-exactness} gives the semantic projection invariant and the
counterexample to cotangent-only completeness.  The algorithmic and complexity
consequences are audited in \cref{sec:algorithm,sec:complexity}.
\Cref{sec:etr-reduction,sec:collapse} state the exact reduction and collapse
obligations.  Boundary behavior, related work, and possible repair directions
appear in the remaining sections.
